% Vietnamese translation for OI Public Library
% Source-PDF: IOI中国国家候选队论文/2005/胡伟栋/附件/polygon.pdf
\documentclass[11pt]{article}
\usepackage{oipl-vn}

\title{IOI'04: Lời giải bài Polygon}
\author{Ioannis Emiris và Elias Tsigaridas}
\date{Ngày 9 tháng 9 năm 2004}

\begin{document}
\maketitle

\origtitle{IOI'04: Solution of Polygon}
\sourcepath{IOI Candidate Team Papers/2005/Hu Weidong/appendix/polygon.pdf}

\section{Phân rã Minkowski}

\subsection{Bối cảnh}

Mặc dù việc tính tổng Minkowski là trực tiếp \cite{guibas-seidel}, việc quyết
định một đa giác có phải là tổng Minkowski của hai đa giác hay không lại là
bài toán NP-đầy đủ.

Với bài toán sau, tồn tại một thuật toán giả đa thức \cite{gao-lauder}; dưới
đây chúng tôi phác thảo thuật toán đó. Nếu dữ liệu vào là dãy \(N\) cạnh, thì
kích thước dữ liệu vào là \(O(N(\log m+\log E))\), trong đó \(m\) là số điểm
nguyên lớn nhất trên một cạnh bất kỳ, còn \(E\) là giá trị tuyệt đối lớn nhất
của một tọa độ dùng để xác định một vectơ cạnh nguyên thủy, như định nghĩa bên
dưới. Thuật toán trong \cite{gao-lauder} chạy trong thời gian
\(O(tNm)\), với \(t=|P\cap\mathbb Z^2|\).

\subsection{Lời giải bài Polygon của IOI-04}

Gọi đa giác vào là \(P\). Giả sử \(P\) gồm \(N\) đỉnh có tọa độ không âm dạng
\(v_i=(x_i,y_i)\), \(0\le i\le N-1\), được cho theo thứ tự ngược chiều kim
đồng hồ.

\paragraph{Định nghĩa 1.1.}
Ta gọi một vectơ \(v=(a,b)\) là \term{nguyên thủy} nếu
\(\gcd(a,b)=1\), với \(a,b\in\mathbb Z_{\ge 0}\). Tương đương, \(v\) nguyên
thủy khi và chỉ khi không có điểm nguyên nào nằm trong phần trong của nó.

Các cạnh của \(P\) được biểu diễn bởi các vectơ
\[
  E_i=v_i-v_{i-1}=(a_i,b_i),\qquad 1\le i\le N,
\]
trong đó \(a_i,b_i\in\mathbb Z\) và các chỉ số được lấy theo modulo \(N\).
Với một cạnh \((a_i,b_i)\), nếu \(n_i=\gcd(a_i,b_i)\), ta xét cạnh nguyên thủy
tương ứng
\[
  e_i=\left(\frac{a_i}{n_i},\frac{b_i}{n_i}\right).
\]
Ta gọi dãy vectơ
\[
  \{E_i\}_{1\le i\le N}=\{n_i e_i\}_{1\le i\le N}
\]
là \term{dãy cạnh} của đa giác, trong đó \(e_i\) là một vectơ nguyên thủy.

Có một quan sát then chốt khi tính các hạng \(A,B\) sao cho \(P=A+B\).

\paragraph{Bổ đề 1.2.}
Mọi cạnh nguyên thủy của \(P\) phải xuất hiện như một cạnh của \(A\) hoặc của
\(B\). Nếu cạnh đó không nguyên thủy, còn có khả năng thứ ba: nó là tổng
Minkowski của một cạnh của \(A\) và một cạnh của \(B\), trong đó hai cạnh này
song song.

Từ đó có thể suy ra sự kiện sau: một đa giác là một hạng của \(P\) khi và chỉ
khi dãy cạnh của nó có dạng
\[
  \{k_j e_j\}_{j\in J},
\]
trong đó \(J\subseteq\{1,\ldots,N\}\), \(0\le k_j\le n_j\),
\(k_j\in\mathbb Z\), và
\[
  \sum_{j\in J} k_j e_j=(0,0),
\]
tức là tổng các vectơ biểu diễn các cạnh của nó bằng không.

Dưới đây chúng tôi trình bày các thuật toán để tìm các hạng theo yêu cầu của
bài Polygon. Khi tính các đa giác hạng, thuật toán được chọn có thể sinh ra
một hạng có một hoặc nhiều tọa độ âm. Trong trường hợp này, ta phải tịnh tiến
cả hai hạng về tọa độ không âm. Đây là một thao tác hợp lệ, vì tịnh tiến không
làm thay đổi bản chất đa giác.

\subsection{Tìm một hạng là đoạn thẳng}

Đa giác có một hạng là đoạn thẳng khi và chỉ khi ít nhất một cặp vectơ con
trong dãy cạnh của nó có tổng bằng vectơ không. Dựa trên sự kiện này, ta trình
bày ba thuật toán với tốc độ tăng dần.

Các thuật toán bắt đầu bằng việc tìm các vectơ nguyên thủy cho mọi cạnh, thông
qua một phép tính \(\gcd\) (mục 1.6 đưa ra hai thuật toán tính \(\gcd\)). Với
những đa giác chỉ chứa các cạnh nguyên thủy, bước này là không cần thiết.

\paragraph{Ngây thơ.}
Trước hết ta tính các cạnh của đa giác. Mỗi cạnh có dạng
\[
  E_i=(x_{i+1}-x_i,\ y_{i+1}-y_i)=(a_i,b_i),
\]
trong đó các chỉ số được lấy theo modulo \(N\). Tiếp theo ta tính tất cả các
giá trị \(\gcd(a_i,b_i)\) và tạo dãy cạnh bằng các vectơ nguyên thủy như ở phần
mở đầu. Việc tính toán này có thể thực hiện trong thời gian \(O(mN)\), với
\[
  m=\max n_i=\max \gcd(a_i,b_i).
\]
Với mỗi vectơ trong dãy, ta tính tổng của nó với mọi vectơ khác trong dãy.
Tổng thời gian là \(O(m^2N^2)\).

\paragraph{Khéo hơn.}
Ta tính dãy cạnh như trước trong thời gian \(O(mN)\). Ta sắp xếp dãy này theo
tọa độ \(x\), trong thời gian \(O(mN\lg(mN))\).

Với mỗi vectơ trong dãy, ta kiểm tra xem có vectơ khác trong dãy với tọa độ
\(x\) đối dấu sao cho tổng của chúng bằng không hay không. Việc tìm kiếm được
thực hiện bằng tìm kiếm nhị phân, vì dãy đã được sắp xếp, trong thời gian
\(O(\lg(mN))\). Tổng thời gian là \(O(mN\lg(mN))\).

\paragraph{Nâng cao.}
Ta tính dãy cạnh như trước trong thời gian \(O(mN)\). Ta chèn các cạnh vào một
bảng băm, dùng tọa độ \(x\) làm khóa. Mỗi phép chèn thực hiện trong \(O(1)\).

Với mỗi cạnh trong dãy, ta tìm trong bảng băm một cạnh khác có tọa độ \(x\) đối
dấu và cộng với nó thành vectơ không. Phép tìm kiếm thực hiện trong \(O(1)\).
Tổng thời gian là \(O(mN)\).

\subsection{Tìm một hạng là tam giác}

\paragraph{Ngây thơ.}
Ta tính dãy cạnh như trước trong thời gian \(O(mN)\). Ta có thể lập mọi bộ ba
có thể và kiểm tra xem chúng có tổng bằng không hay không. Tổng thời gian là
\(O(m^3N^3)\).

\paragraph{Khéo hơn.}
Ta tính dãy cạnh như trước trong thời gian \(O(mN)\) và sắp xếp nó theo tọa độ
\(x\) trong thời gian \(O(mN\lg(mN))\).

Ta lập mọi tổng có thể của hai vectơ trong thời gian \(O(m^2N^2)\). Với mỗi
tổng như vậy, ta dùng tìm kiếm nhị phân trong thời gian \(O(\lg(mN))\) để tìm
một vectơ khác sao cho tổng của chúng bằng không. Tổng thời gian là
\(O(m^2N^2\lg(mN))\).

\paragraph{Nâng cao.}
Ta tính dãy cạnh như trước trong thời gian \(O(mN)\) và sắp xếp chúng theo thứ
tự tăng dần của tọa độ \(x\) trong thời gian \(O(mN\lg(mN))\).

Với mỗi cạnh trong dãy, gọi là \(k\), quét dãy từ trái để tìm \(i\) và từ phải
để tìm \(j\), sao cho
\[
  E_k+E_i+E_j=0.
\]
Nhờ thứ tự đã sắp xếp, ta có thể tăng \(i\) hoặc giảm \(j\) tùy theo tổng
\(E_k+E_i+E_j\) là dương hay âm. Tổng thời gian là \(O(m^2N^2)\).

Một cách khác là chèn các vectơ vào bảng băm, dùng tọa độ \(x\) làm khóa, và
với mỗi tổng của hai vectơ, ta tìm trong bảng băm một vectơ trong thời gian
\(O(1)\) sao cho tổng của cả ba bằng không. Khi đó tổng thời gian cũng là
\(O(m^2N^2)\).

\subsection{Tìm một hạng là tứ giác}

\paragraph{Ngây thơ.}
Ta tính dãy cạnh như trước trong thời gian \(O(mN)\). Ta có thể lập mọi bộ bốn
có thể và kiểm tra xem chúng có tổng bằng không hay không. Tổng thời gian của
thuật toán là \(O(m^4N^4)\).

\paragraph{Khéo hơn.}
Ta tính dãy cạnh như trước trong thời gian \(O(mN)\). Ta có thể lập mọi bộ ba
trong thời gian \(O(m^3N^3)\). Với mỗi bộ ba, ta tìm một vectơ sao cho tổng của
cả bốn bằng không.

Nếu đã sắp xếp các cạnh theo tọa độ \(x\) và dùng tìm kiếm nhị phân, tổng thời
gian của thuật toán là \(O(m^3N^3\lg(mN))\).

Nếu chèn các cạnh vào bảng băm, dùng tọa độ \(x\) của chúng làm khóa, ta thực
hiện phép tìm kiếm trong thời gian \(O(1)\), do đó tổng thời gian của thuật
toán là \(O(m^3N^3)\).

\paragraph{Nâng cao.}
Ta tính dãy cạnh như trước trong thời gian \(O(mN)\). Ta lập mọi tổng có thể
của hai cạnh trong thời gian \(O(m^2N^2)\) và chèn chúng vào một bảng băm, dùng
tọa độ \(x\) làm khóa. Với mỗi tổng như vậy, ta tìm trong bảng băm một vectơ
sao cho tổng toàn bộ bằng không. Tổng thời gian là \(O(m^2N^2)\).

\subsection{Tính \texorpdfstring{\(\gcd\)}{gcd}}

Chúng tôi trình bày hai thuật toán tính \(\gcd\) của hai số nguyên \(a,b\).
Thuật toán thứ nhất là thuật toán Euclid truyền thống, còn thuật toán thứ hai
là thuật toán \(\gcd\) nhị phân (xem chương 4 của \cite{knuth}).

\paragraph{Thuật toán 1. Euclid \(\gcd(a,b)\).}
\begin{description}[leftmargin=7em,style=nextline]
  \item[Yêu cầu:] \(a,b\in\mathbb Z\).
  \item[Bảo đảm:] \(y=\gcd(a,b)\).
\end{description}
\begin{verbatim}
if b = 0 then
  RETURN a
else
  RETURN Euclid gcd(b, a mod b)
end if
\end{verbatim}

\paragraph{Thuật toán 2. \(\gcd\) nhị phân \(\gcd(a,b)\).}
\begin{description}[leftmargin=7em,style=nextline]
  \item[Yêu cầu:] \(a,b\in\mathbb Z\).
  \item[Bảo đảm:] \(y=\gcd(a,b)\).
\end{description}
\begin{verbatim}
g = 1
while a is even AND b is even do
  a = a/2 (right shift)
  b = b/2
  g = 2 * g (left shift)
end while
while u > 0 do
  if a is even then
    a = a/2
  else if b is even then
    b = b/2
  else
    t = |a - b|/2
  end if
  if a < b then
    b = t
  else
    a = t
  end if
end while
RETURN g * b
\end{verbatim}

\begin{thebibliography}{9}
\bibitem{gao-lauder}
S. Gao và A. Lauder,
\textit{Decomposition of polytopes and Polynomials},
Discrete and Computational Geometry 26 (2001), 501--520.

\bibitem{guibas-seidel}
L. J. Guibas và R. Seidel,
\textit{Computing Convolutions by Reciprocal Search},
Symposium of Computational Geometry, 2001.

\bibitem{knuth}
D. Knuth,
\textit{The Art of Computer Programming}, tập 2,
Addison-Wesley, ấn bản thứ 3, 1998.
\end{thebibliography}

\end{document}
